% preamble.tex

% =========================
% Idioma y codificación
% =========================
\usepackage[spanish]{babel}
\usepackage[utf8]{inputenc}
\usepackage[T1]{fontenc}

% =========================
% Fuente
% =========================
\usepackage{newtxtext,newtxmath}

% =========================
% Ajuste fino para cajas y diagramas
% =========================
\usepackage{adjustbox}

% =========================
% Márgenes
% =========================
\usepackage{geometry}
\geometry{
  left=2.54cm,
  right=2.54cm,
  top=2.54cm,
  bottom=2.54cm
}

% =========================
% Interlineado
% =========================
\usepackage{setspace}
\doublespacing

% =========================
% Formato de párrafo
% =========================
\setlength{\parindent}{1.27cm}
\setlength{\parskip}{0pt}

% =========================
% Mejoras tipográficas
% =========================
\usepackage{microtype}
\emergencystretch=3em
\hfuzz=6pt
\vbadness=10000
\hbadness=10000

% =========================
% Figuras y tablas
% =========================
\usepackage{graphicx}
\usepackage{caption}
\captionsetup{
  font=small,
  labelfont=bf
}

\usepackage{booktabs}
\usepackage{array}
\usepackage{tabularx}
\usepackage{longtable}
\usepackage{ltablex}
\keepXColumns
\usepackage{threeparttablex}

\usepackage{float}
\usepackage{placeins}
\usepackage{enumitem}

% =========================
% Columnas útiles para tablas
% =========================
\newcolumntype{Y}{>{\raggedright\arraybackslash}X}
\newcolumntype{L}[1]{>{\raggedright\arraybackslash}p{#1}}

% =========================
% Configuración APA-friendly para tablas
% =========================
\newcommand{\APAtablesetup}{%
  \small
  \setlength{\tabcolsep}{4pt}%
  \renewcommand{\arraystretch}{1.2}%
}

\newcommand{\APAtablecontinued}{\footnotesize\textit{Continúa en la siguiente página.}}

\newenvironment{APAlongtable}[3]{%
  \APAtablesetup
  \begin{center}
  \begin{longtable}{#1}
  \caption{#2}\label{#3}\\
  \toprule
}{%
  \bottomrule
  \end{longtable}
  \end{center}
}

% =========================
% Hipervínculos y referencias
% =========================
\usepackage{hyperref}
\hypersetup{
  colorlinks=true,
  linkcolor=black,
  citecolor=black,
  urlcolor=black,
  pdftitle={Tesis doctoral},
  pdfauthor={José Martínez}
}

\usepackage[nameinlink,noabbrev]{cleveref}
\crefname{figure}{Figura}{Figuras}
\crefname{table}{Tabla}{Tablas}
\crefname{chapter}{Capítulo}{Capítulos}
\crefname{section}{Sección}{Secciones}

% =========================
% Bibliografía
% =========================
\usepackage{csquotes}
\usepackage[
  backend=biber,
  style=apa,
  citestyle=apa,
  sorting=nyt,
  maxcitenames=2,
  maxbibnames=99,
  uniquelist=false
]{biblatex}

\DeclareLanguageMapping{spanish}{spanish-apa}

% Base global
\addbibresource{references/global.bib}

% Bases por capítulo
\addbibresource{references/ch01.bib}
\addbibresource{references/ch02.bib}
\addbibresource{references/ch03.bib}
\addbibresource{references/ch04.bib}
\addbibresource{references/ch05.bib}
\addbibresource{references/ch06.bib}
\addbibresource{references/ch07.bib}

% =========================
% TikZ
% =========================
\usepackage{tikz}
\usetikzlibrary{
  positioning,
  arrows.meta,
  shapes.geometric,
  calc,
  fit
}

% =========================
% Ajustes para evitar conflictos típicos con babel y TikZ
% =========================
\AtBeginDocument{%
  \shorthandoff{<>}
}

% =========================
% Comandos consistentes
% =========================
\newcommand{\DUA}{Diseño Universal para el Aprendizaje (DUA)}
\newcommand{\DBR}{Design-Based Research (DBR)}
\newcommand{\IA}{Inteligencia Artificial (IA)}

% =========================
% Mejor comportamiento general
% =========================
\raggedbottom


\usepackage{caption}
\captionsetup{
  font=small,
  labelfont=bf,
  skip=8pt
}

\usepackage{pdflscape}